% Authority: exact-scope leaf 5; full PDF page 84; printed p.67.
% U:p067.u01 | paragraph-start
\par\indent
“If the numbers $A$ and $A'$ are both biquadratic residues or both biquadratic
non-residues with respect to $p$, the product $AA'$ will be a biquadratic residue
with respect to $p$; if, on the contrary, one of the numbers $A$ and $A'$ is a
residue and the other a biquadratic non-residue relative to $p$, $AA'$ will be a
biquadratic non-residue with respect to $p$.”
% U:p067.u02 | paragraph-start
\par\indent
Taking $A'=-1$, and having regard to what precedes, one will see that, if $p$
is of the form $8n+1$, $A$ will at the same time as $-A$ be a biquadratic residue
or non-residue, and that, on the contrary, if $p$ is of the form $8n+5$, one of
the numbers $A$ and $-A$ will be a residue and the other a biquadratic non-residue
with respect to $p$.
% U:p067.section | heading
\sectionnumber{2}
% U:p067.u03 | paragraph-start; continues through p.68 and p.69
\par\indent
After establishing these preliminaries, we shall deal with the search for the
characters which distinguish the prime divisors of the formula $x^4-2$. It is
known that the prime divisors of $x^2-2$ are of one of these two forms:
$8n+1$, $8n+7$, and that reciprocally every prime number of one of these forms
divides the formula $x^2-2$. According to what we said in the preceding paragraph,
we need not take account of the latter of these two forms, and it will suffice
to consider the prime numbers $8n+1$. Let $p$ be a prime number of this kind,
and put, as one is permitted to do, $p=t^2+2u^2$, where $u$ will be even and
$t$ odd. Let $u=2^\nu kk'k''\ldots$, $2^\nu$ being the highest power of $2$ which
divides $u$, and $k$, $k'$, $k''$, $\ldots$ designating odd prime numbers, several
of which may be equal to one another. The equation $t^2+2u^2=p$ immediately
gives $t^2\equiv p\md{k}$. The number $p$ is therefore a quadratic residue
with respect to $k$, which we shall write thus: $\leg{p}{k}=1$, adopting the
notation used by Mr. \textsc{Legendre}. Recall that, if $c$ designates a prime
number and $M$ any number not divisible by $c$, this illustrious geometer uses
the sign $\leg{M}{c}$ to designate the remainder which one obtains by dividing
by $c$ the power $M^{\frac{c-1}{2}}$, a remainder which is known to be equal to
$1$ or $-1$, according as $M$ is or is not a quadratic residue with respect to $c$.
Applying to the relation $\leg{p}{k}=1$ the theorem known under the name law
of reciprocity (the \emph{theorema fundamentale} of Mr. \textsc{Gauss}), one will have, $p$ being
